\section{Venskabsforeninger}

\TKET har igennem sine mange år fået sig et par venskabsforeninger, som er beskrevet her.


\subsection{Nabla}
Over årende har vi haft besøg af og besøgt foreningen \textit{Nabla}. Som der står på deres hjemmeside er Nabla: ``... linjeforeningen for studenter ved sivilingeniørstudiet Fysikk og Matematikk ved NTNU, ..."\footnote{https://nabla.no/}, i Trondheim, Norge. Nabla kontaktede \TKET omkring midt 10'erne og inviterede en håndfuld medlemmer til Norge, og siden da har \TKET sendt medlemmer til besøg i Norge, og de ligeledes. Efter corona-tiden er hyppigheden af disse turen desværre dalet. Sidste besøg i Trondheim var i 2023, da \TKET blev inviteret til Nablas 80 års jubilæum.

Disse ture har skabt en rig kulturudveksling, blandt andet blev Lambo-sangen bragt tilbage til Danmark efter et besøg hos Nabla. Da man i Norge ikke må drikke i offentligheden, bruger de selv sangen til at tømme deres drikke, inden de tager fra forfest videre til hovedfesten. Her synges Lambo-sangen, flasken sendes rundt i en cirkel, og når man får flasken drikke man en lille smule, indtil flasken er tom. Da alkohol er meget billigere i Danmark end i Norge, har den så udviklet sig til at være en form for bunde-cirkel. Her bunder folk en øl og synger slutningen, inden man går til den næste der skal bunde.

\begin{center}
	\textit{Berseblæsten - Lambo Glück Auf (1978)}
	\parbox[t]{.5\textwidth}{%
	\begin{itemize}
		\item[] \textit{Tilskuere synger}
		\setlength{\itemsep}{-5pt}
		\item[1.] Se der står en fyllehund, mine herrer lambo.
		\item[] Sett nu flasken for din munn, mine herrer lambo.
		\item[] Se hvordan den dråpen vanker ned ad halsen på den dranker
		\item[] Lambo, lambo, mine herrer lambo \textit{(gentages til vedkommende har drukket op)}
	\end{itemize}
}%
	\parbox[t]{.5\textwidth}{%
		\begin{itemize}
			\setlength{\itemsep}{-5pt}
			\item[] \textit{Drikkeren synger}
			\item[2.] Jeg mitt glass utdrukket har, mine herrer lambo.
			\item[] Se der fins ei dråpen kvar, mine herrer lambo.
			\item[] Som bevis der på jeg vender, flasken på dens rette ende \textit{(drikkeren vender flasken på toppen af hovedet som bevis)}
		\end{itemize}
	}%
\end{center}
\vspace{-20pt}
	\parbox[t]{.5\textwidth}{%
	\begin{itemize}
		\item[] \textit{Tilskuere synger}
		\setlength{\itemsep}{-5pt}
		\item[3.] Lambo, lambo, mine herrer lambo
		\item[] Hun/han kunne kunsten hun/han var et jævla fyllesvin.
		\item[] Så går vi til nestemann og ser hva hun/han formår. \textit{(Find en ny person, og forsæt til man ikke gider mere)}
	\end{itemize}
}%

\newpage

\subsection{Caféen?}
\citatliste[r]{0.46\linewidth}{%
    ``TK er en kult der tilbeder en overfermenteret ost.''\\[6pt]
    ``Deres ritualer inkluderer, men er ikke begrænset til: nøgenkost, uddeling af flyers, bygning af tømmerflåder og livestreaming af opkast.''\\[6pt]
    ``Jeg elsker dem for alt de er og gør for studiemiljøet.''%
}

På Københavns Universitet ligger der, et sted på campus, en bygning, som i daglig tale kaldes ``Caféen?'' eller blot ``C?''. På Caféen? kan man i løbet af dagen få sig en kop kaffe og sidde ved hyggelige borde og sludre. Men C? bliver nu mere anset for at være en bar end en café, hvilket den godt 6 meter lange bardisk i lokalet også vidner om.
Foreningen blev stiftet i 1993, og er passende i vedtægterne navngivet ``Foreningen til Cafédrift af 1993'', selvom det dog er et tungt navn at bruge i hverdagssnak. Caféen? dækker over studier som datalogi, matematik, fysik, kemi, og diverse andre tilhørende studier.
På mange måder kan Caféen? anses som KU's bud på \TKET.

Den observante læser har nu allerede spottet at Caféen? har lidt problemer med tegnsætningen, og det er ganske bevidst fra Caféen?s side af. Spørgsmålstegnet er en del af navnet, da det blot at nævne navnet byder op til spørgsmålet, om man ikke skulle se at komme ned på Caféen?.
C? bliver også af flere kaldet ``hjem'', hvilket både refererer til den mørke tid mellem midt 2004 og 2007 hvor C? manglede lokaler og derfor var hjemløse, men også til et Basserne vægmaleri i lokalet, hvor S'janten og Plato sidder i en bunke af tomme øldåser, og diskuterer, om de ikke skulle se at komme hjem, blot for at konkludere, at de allerede er hjemme.

På \TKET lærte vi for nyligt C? at kende i BEST 22/33, da de rakte ud til os for at høre om de kunne være med i kapsejladsen. Efter at måtte melde til dem, at det nu ikke var noget vi kunne hjælpe med, fandt vi ud af, at vi begge havde hængere, og at de begge kunne lide øl, så der skulle naturligvis sættes en legeaftale op.

\citatliste[l]{0.46\linewidth}{%
    ``Snakker vi alvorligt hvad man tror de laver eller sjove udtalelser?''\\[6pt]
    ``Also yes, TK er cult. Men de for dumme til at kunne skrive Sigma tegnet.''%
}

Og hvilken bedre mulighed end Hej-dag. Op dukkede vores hængere, hvor de først skulle lege gætteleg, om hvad dagen mon skulle byde på. De havde lige nået til at tro, at der faktisk ikke var planlagt noget, indtil der dukkede en bus op på parkeringspladsen fyldt med københavnere. Ligeså havde C? fortalt deres hængere, at de skulle til BonBon land, og ikke den lange vej rundt over Storebælt.

\citat[r]{0.46\linewidth}{Jeg har hørt rygter om at \TKET ligger lige ved BonBon Land. Samme bus kører ihvertfald begge steder hen.}

Men her stod de. Der blev sunget velkommen, udvekslet porrer, og C? fik sig en kammerintro. Vi holdte fælles frokost, efterfulgt af en rask løvfaldstur i uniparken. Herefter stod den på fri leg, hvor der blev delt diverse spil og traditioner på tværs. Efter en gang pizza, var vi klar til torsdagsmødet, som trak sig så langt ud på natten, som kasserne på bordet tillod sig. Alt i alt, en god legeaftale.

En af de spil C? havde med var ``racerbile'', som spilles på en malet bane af hexagoner. Det gør det en smule sværre at bare spille, men heldigvis fik vi i gave en plade til \TKET. På C? er det en smule nemmere, da flere af bordene passende allerede har malet en bane på.
I spillet racer man rundt på hexagonerne med hver sin kapsel, som var det en racerbane. Der følger med spillet en lang række regler, men da en af dem er, at man skal drikke når man forklarer regler, vil jeg for mit eget velbefindende stoppe her\footnote{Jeg har drukket for den her.}, og blot bemærke, at det er et ganske godt spil at fordrive tiden med i godt selvskab.

Ligesom vi på \TKET har Ceres Top som stam-øl, har de på C? Tuborg Gulddamen, som de nyder i store drag.
Faktisk så store drag, at der på C? kun er 11 tåre i en gulddame på 33cl, hvilket giver 8 tåre i en almindelig øl. Man skal derfor passe på med at tage kammerspil med på Caféen?, da man helt automatisk bliver $1,375$ gange så fuld som normalt. Til academy har de dog heldigvis adopteret den standartiserede 14 tåres øl.

Så hvordan finder jeg så dette magiske sted? Tag til København, find KU, stavrer rundt på campus efter lugten af øl, til du finder en lille bygning med en rød dør. Træd indenfor, køb dig et medlemsskab og en øl, og slå dig en sludder af til den lyse morgen.

%Så.. Caféen??

\subsection{F-Klubben}
Det siges, at Niels Wendell Pedersen (NWP) satte sig på en ølkasse midt ude på mark og sang, hvorved Aalborg Universitets campus opstod rundt omkring ham.
Der opstod i NWPs forkontor, hvor F-sektorens sekretær, Hanne, arbejdede, en kaffemaskine, en sofa og et par stole. I denne kaffestue var der også et køleskab og en kasse med sedler til at holde styr på køb. F-Klubben blev her stiftet 31. februar 1976. 
Tiden går, uddannelserne ændrer sig, F-sektoren eksisterer ej længere. Såvel F-Klubben som \TKET har jubilæum i år, F-Klubben 50 år og \TKET 70 år. Til dette giver F-Klubben et stort TILLYKKE og en endnu større SKÅL!

Men nok om dette. I forbindelse med jubilæet, benådes I herved F-Klubbens ældgamle regler i den endnu ældre disciplin- ØL-LUDO, også kendt som F-LUDO.

\subsubsection{Dommer og sidedommer}
For at ØL-LUDO kan blive afviklet på korrekt vis, har dommeren, som er streng, dybt uretfærdig og absolut modtagelig over for bestikkelse, sin sidedommer. Denne skal være dommeren behjælpelig med kontrol af overholdelse af diverse regler. Det teoretiseres, at dommer og sidedommer mangler briller og/eller er natteblinde, men det kan ikke bekræftes, da ØL-LUDO selvfølgelig foregår i lyset.

\subsubsection{Regler - en nødvendighed}
Hvert hold, som består af 2-3 deltagere (2 drenge, 3 piger, 2 piger og 1 dreng eller 1 dreng og 1 pige), tildeles en kasse øl for egen regning, og kan efter behov tildeles flere. Spillet startes derefter med, at alle hold samtidigt bunder en øl, og det hold, der bliver først færdig, starter.

\parbox[t]{.5\textwidth}{%
	\begin{itemize}
	\item Er alle 4 brikker i fængsel, har holdet 3 forsøg til at slå en sekser og derved få en brik ud.
	\item Slag med den hellige F-erning skal ske med respekt og dyd. Slag udføres aldrig i vrede, og F-erningen må hverken kastes, smides eller tiltales i nedrig tone. Vanhelligelse sones efter dommernes idømmelse og bødlens faglighed.
	\item Ved slag af sekser opnås til enhver tid ekstra slag.
	\item Ved slag af flogo rykkes én brik frem til den første umiddelbart efterfølgende F-logo.
	\item Ved slag af stjerne rykkes en brik frem til den første umiddelbart efterfølgende stjerne.
	\item Kan et slag ikke udnyttes, drikker holdet i stedet en øl. Dette regnes ikke som en straf, men som trøst.
	\item Der er helle på alle F-logoer, dog kan et hold som slår en brik ud af fængsel og derved lander på egen F-logo slå en brik i fængsel.
	\item Det giver helle at have to brikker på samme felt.
	\item Lander man på en stjerne, uden at have slået en stjerne, springer man videre til den umiddelbart efterfølgende stjerne. Hvis den første stjerne er optaget, slås denne brik i fængsel, og man bliver stående. Hvis ikke dette er tilfældet, men der står en brik på den anden stjerne, rykker brikken selv i fængsel.

\end{itemize}
	}%
	\parbox[t]{.5\textwidth}{%
		\begin{itemize}
	\item Kan man slå en brik i fængsel, skal man gøre dette. Hvis dette undlades, bestemmer dommeren, hvilken af holdets egne brikker der skal i fængsel.
	\item Slås en brik i fængsel, eller rykker man en brik hjem er straffen én øl.
	\item Der må kun være 4 af hvert holds brikker på ØL-LUDO-brættet.
	\item Alle brikker skal være identificerbare som tilhørende et af holdene.
	\item Brikkerne skal være fulde.
	\item Et hold kan til enhver tid idømmes straf-roulader, -chips, -øl mm. for utilbørlig opførsel eller mangel på samme.
	\item Under strafafsoningen kan holdet ikke deltage i det videre ØL-LUDO-spil, og sidder derfor over.
	\item Når en straf er afsonet, holdes flasken over hovedet med bunden i vejret. Dette skal kontrolleres af enten sidedommeren eller dommeren selv.
	\item Når en brik er nået op på opløbsstrækningen, skal den flyttes direkte hjem. Dette er ensbetydende med, at stjerne svarer til 2 og F-logo til 5.
	\item Det hold, der først får alle 4 brikker hjem, har vundet, og spillet stoppes umiddelbart herefter.
	\item Det er ikke tilladt at gå på ØL-LUDO-brættet eller på anden måde at optræde ærekrænkende over for ØL-LUDO-ånden, heller ikke i tankerne.
\end{itemize}
}%
\begin{center}
		\textbullet\ \ Dommeren har altid ret.
\end{center}
